\section{Introduction}
The Curry-Howard isomorphism shows the deep relationship between intuitionistic logic and computation. Specifically, it establishes that proofs can be interpreted as programs and proof reduction as computation. Because of this, we can see that the different formalizations of intuitionistic logic will result in different computational interpretations. Howard, for example, established the connection between natural deduction and the simply-typed $\lambda$-calculus. Curry took a different approach, associating Hilbert-style proofs to combinators. Herbelin established the system LJT based on the sequent calculus with a stoup, leading to a computational interpretation featuring explicit substitutions. 

Adding to this lineage, DeYoung et al.~\cite{DeYoung2020} presents the semi-axiomatic sequent calculus (SAX), a formalization of intuitionistic logic which arises from by replacing all non-invertible rules in the sequent calculus with axioms.

We note that the rules that are maintaned from the standard sequent calculus (invertible rules) are the right rules from negative connectives ($A \supset B$, $A \& B$), and the left rules of positive connectives ($A \otimes B$, $\textbf{1}$, $A \lor B$)~\cite{DeYoung2020}. These rules can be seen as asynchronous in that they can be applied eagerly in bottom-up proof construction. Conversely, the left rules of negative connectives and the right rules of positive connectives may be postponed~\cite{DeYoung2020}. These properties, along with t

\begin{figure}[ht!]
\[
  \infer[\mathsf{Id}_A]{\Gamma, A \vdash A}{}
  \qquad
  \infer[\mathsf{Cut}_A]{\Gamma \vdash C}{\Gamma \vdash A & \Gamma, A \vdash C}
\]
\[
  \infer[{\supset}R]{\Gamma \vdash A \supset B}{\Gamma, A \vdash B}
  \qquad
  \infer[{\supset}L^0]{\Gamma, A, A \supset B \vdash B}{}
\]
\[
  \infer[{\otimes}R^0]{\Gamma, A, B \vdash A \otimes B}{}
  \qquad
  \infer[{\otimes}L]{\Gamma, A \otimes B \vdash C}{\Gamma, A \otimes B, A, B \vdash C}
\]
\[
  \infer[\mathbf{1}R^0]{\Gamma \vdash \mathbf{1}}{}
  \qquad
  \infer[\mathbf{1}L]{\Gamma, \mathbf{1} \vdash C}{\Gamma, \mathbf{1} \vdash C}
\]
\[
  \infer[{\with}R]{\Gamma \vdash A \mathbin{\with} B}{\Gamma \vdash A & \Gamma \vdash B}
  \qquad
  \infer[{\with}L^0_1]{\Gamma, A \mathbin{\with} B \vdash A}{}
  \qquad
  \infer[{\with}L^0_1]{\Gamma, A \mathbin{\with} B \vdash B}{}
\]
\[
  \infer[{\lor}R^0_1]{\Gamma, A \vdash A \lor B}{}
  \qquad
  \infer[{\lor}R^0_2]{\Gamma, B \vdash A \lor B}{}
  \qquad
  \infer[{\lor}L]{\Gamma, A \lor B \vdash C}{\Gamma, A \lor B, A \vdash C & \Gamma, A \lor B, B \vdash C}
\]
\caption{The Semi-Axiomatic Sequent Calculus (\textsc{Sax})~\cite{DeYoung2020}.}
\label{fig:sax-rules}
\end{figure}


